\subsection{章末確認問題}

\begin{enumerate}
  \item ソフトウェア設計という語が持つ四つの意味を説明せよ。
  \item 設計思考の五段階を、自分の言葉で説明せよ。
  \item 高水準設計と詳細設計の違いを、外部向き・内部向きという観点から説明せよ。
  \item 抽象化、モジュール化、カプセル化の違いを説明せよ。
  \item 結合度を低くし、凝集度を高くすることが望ましい理由を説明せよ。
  \item 並行性を設計するとき、同期や排他制御が必要になる理由を説明せよ。
  \item 構造設計記述と振る舞い設計記述の例を、それぞれ三つ挙げよ。
  \item 設計パターンが「語彙」として役立つ理由を説明せよ。
  \item ドメイン駆動設計で、共有語彙が重要になる理由を説明せよ。
  \item 設計レビュー、静的分析、プロトタイピングはどのように使い分けるか。
\end{enumerate}

\begin{pointbox}{解答の方向性}
1は「分野、プロセス、結果、ライフサイクル段階」である。3は、高水準設計が外部とのやり取りや主要構成要素の責務を扱い、詳細設計が内部モジュール、データ構造、アルゴリズムを扱う点を説明する。5は、変更影響の局所化と理解容易性を軸に説明できる。7は、構造ならクラス図、コンポーネント図、ERD、配置図など、振る舞いなら活動図、シーケンス図、DFD、状態遷移図、決定表などを挙げればよい。
\end{pointbox}
